\begin{myexercise}
    In dieser Aufgabe erweitern Sie den Code aus der Aufgabe \textbf{"ZweiDFormen"}, um eine Klassenhierarchie für geometrische Körper unterschiedlicher Dimensionen zu schaffen.

    \subsection*{1. Die Basisklasse schaffen}
    \begin{enumerate}
        \item Erstellen Sie die \textbf{Basisklasse} \texttt{GeometrischerKoerper}.
        \begin{itemize}
            \item Setzen Sie das \textbf{Klassenattribut} \texttt{dimension} auf \texttt{None}.
            \item Der \textbf{Konstruktor} \texttt{\_\_init\_\_(self, name)} muss ein \textbf{Instanzattribut \texttt{self.name}} speichern.
        \end{itemize}

        \item Passen Sie die Klasse \texttt{ZweiDForm} an: Sie soll von \textbf{\texttt{GeometrischerKoerper} erben}.
        \item Setzen Sie in \texttt{ZweiDForm} das Attribut \texttt{dimension} auf den Wert \texttt{2}.
        \item Testen Sie die Vererbung mit einem 2D-Objekt.
    \end{enumerate}

    \subsection*{2. Erweiterung auf 3D}
    \begin{enumerate}
        \setcounter{enumi}{4} % Zähler fortsetzen
        \item Fügen Sie eine Klasse \texttt{DreiDKoerper} hinzu, die von \textbf{\texttt{GeometrischerKoerper} erbt}.
        \item Setzen Sie in \texttt{DreiDKoerper} das Attribut \texttt{dimension} auf \texttt{3}.
        \item \textbf{Implementieren} Sie in \texttt{DreiDKoerper} die Methoden \texttt{berechne\_oberflaeche(self)} und \texttt{berechne\_volumen(self)}. Diese Methoden sollen in der Basisklasse \texttt{DreiDKoerper} eine Warnung drucken: Die Methode <methodenname> muss in der Unterklasse {self.name} definiert werden.
    \end{enumerate}

    \subsection*{3. Konkrete 3D-Formen}
    \begin{enumerate}
        \setcounter{enumi}{7} 
        \item Erstellen Sie die Klasse \texttt{Wuerfel}, die von \texttt{DreiDKoerper} erbt.
        \begin{itemize}
            \item \textbf{Erstellen} Sie den \textbf{Konstruktor} und \textbf{bestimmen} Sie dabei, welche \textbf{Instanzattribute}  benötigt werden, um die Oberfäche und das Volumen berechnen zu können.
            \item Implementieren Sie \texttt{berechne\_oberflaeche()}: $6 \cdot a^2$, wo $a:=$ die Seitenlänge des Würfels ist.
            \item Implementieren Sie \texttt{berechne\_volumen()}: $a^3$.
        \end{itemize}

        \item Erstellen Sie die Klasse \texttt{Kugel}, die von \texttt{DreiDKoerper} erbt.
        \begin{itemize}
            \item \textbf{Erstellen} Sie den \textbf{Konstruktor} und \textbf{bestimmen} Sie dabei, welche \textbf{Instanzattribute} benötigt werden, um die Oberfäche und das Volumen berechnen zu können.
            \item Implementieren Sie \texttt{berechne\_oberflaeche()}: $4 \cdot \pi \cdot r^2$, wo $r:=$ der Radius der Kugel ist.
            \item Implementieren Sie \texttt{berechne\_volumen()}: $\frac{4}{3} \cdot \pi \cdot r^3$.
        \end{itemize}

        \item \textbf{Testen} Sie die Implementierung mit Objekten von \texttt{Wuerfel} und \texttt{Kugel} (Berechnung von Oberfläche und Volumen).
    \end{enumerate}

\end{myexercise}